\documentclass[11pt]{article}
\usepackage[utf8]{inputenc}
\usepackage[T1]{fontenc}
\usepackage[a4paper,margin=0.78in]{geometry}
\usepackage{graphicx,subcaption,float,booktabs,array,enumitem,microtype,xcolor,listings,amsmath,amssymb,tabularx}
\usepackage[hidelinks]{hyperref}
\setlist{nosep,leftmargin=*}
\setlength{\parindent}{0pt}
\setlength{\parskip}{4pt}
\setlength{\emergencystretch}{2em}
\setcounter{tocdepth}{2}
\renewcommand{\arraystretch}{1.12}
\graphicspath{{plots/}}
\newcommand{\plotone}[4][0.9\linewidth]{%
  \begin{figure}[H]\centering\includegraphics[width=#1]{#2}\caption{#3}\label{#4}\end{figure}}
\newcommand{\plotpair}[6]{%
  \begin{figure}[H]\centering
  \begin{subfigure}[t]{0.48\textwidth}\centering\includegraphics[width=\linewidth]{#1}\caption{#2}\end{subfigure}\hfill
  \begin{subfigure}[t]{0.48\textwidth}\centering\includegraphics[width=\linewidth]{#3}\caption{#4}\end{subfigure}
  \caption{#5}\label{#6}\end{figure}}
\newcommand{\StudentOne}{Abhi Jain (23B0903)}
\newcommand{\StudentTwo}{Deshank Pratap Singh (23B0938)}
\newcommand{\StudentThree}{Dhruv Jain (23B0942)}
\newcommand{\StudentFour}{Sameer Arvind Patil (23B1035)}
\begin{document}
\hypersetup{pageanchor=false}
\pagenumbering{alph}
\begin{titlepage}
    \centering
    {\scshape\bfseries\Huge Advanced Computer Architecture\par}
    \vspace{0.2cm}{\scshape\bfseries\Huge (CS-683)\par}
    \vspace{0.8cm}{\huge\textbf{Programming Assignment 2\\Report}\par}
    \vspace{0.8cm}{\Large Data Prefetching and Cache Replacement Policies\par}
    \vspace{1cm}{\Large\StudentOne\par}{\Large\StudentTwo\par}{\Large\StudentThree\par}{\Large\StudentFour\par}
    \vfill{\large Autumn 2026\par}
    {\large ChampSim measurements and implementation analysis\par}
\end{titlepage}
\hypersetup{pageanchor=true}
\pagenumbering{roman}\tableofcontents\newpage\pagenumbering{arabic}

\section{Overview and Experimental Methodology}
\subsection{Objectives}
This report analyzes the provided IP-stride prefetcher (Task 1A), the implemented Stride Pattern Prefetcher (SPP, Task 1B), instruction-conscious RRIP (Task 2A), dynamic I-RRIP (Task 2B), and a combined SPP/replacement configuration (Task 3). The report separates source-code behavior from measured performance and identifies places where the supplied experiment provenance does not support a stronger conclusion.

\subsection{Experimental setup}
\begin{table}[H]\centering
\caption{Configuration recorded in the experiment runner and build script.}
\begin{tabularx}{\linewidth}{>{\raggedright\arraybackslash}p{0.34\linewidth}X}
\toprule Parameter & Recorded setting \\
\midrule
Simulator & ChampSim repository/fork used by the supplied runner \\
Trace names & \texttt{Tr1.gz}, \texttt{Tr2.gz}, \texttt{Tr3.gz}; semantic workload descriptions were not supplied \\
Warmup & 25,000,000 instructions per run \\
Measured instructions & 25,000,000 instructions per run \\
Core count & 1 \\
Branch predictor & \texttt{hashed\_perceptron} \\
L1D prefetcher & Task-dependent: \texttt{no}, \texttt{ip\_stride}, or \texttt{spp} \\
L2 replacement & Task-dependent: LRU, RRIP, I-RRIP, or dynamic I-RRIP \\
Other cache/TLB prefetchers and replacement & Build script sets prefetchers to \texttt{no} and other replacement policies to LRU unless configured by the task \\
Cache parameters & Default repository settings; no separate cache-configuration override was provided \\
Host/compiler details & Not recorded in the supplied experiment metadata \\
Repetitions & One recorded result per trace/configuration/threshold; no repeated-run variance estimate supplied \\
\bottomrule
\end{tabularx}\end{table}

\subsection{Metrics and baselines}
IPC is the simulator-reported instructions-per-cycle value. Speedup is the configuration IPC divided by the task-specific baseline IPC. Task 1B uses no prefetching + LRU as baseline; Tasks 2A and 2B use RRIP as baseline; Task 3 uses no prefetching + LRU as baseline. The baseline therefore differs across tasks and must not be mixed when comparing reported speedups.

MPKI is misses per thousand measured instructions. L1D and aggregate L2C load MPKI are derived from reported load-miss counters using 25M as the measured-instruction denominator when the parser does not find a direct MPKI line. L2 instruction-load and L2 data-load MPKI are separate simulator counters. In these files, the aggregate L2C load-MPKI field and the instruction/data category counters do not add up to one another; they should be interpreted as separately reported statistics, not as arithmetic components of a common total.

Prefetch accuracy is reported as the percentage of issued L1D prefetches later used, using the runner's parsed useful/issued counters where available. It is not applicable to a no-prefetch configuration, so the tables use N/A rather than interpreting the raw zero as measured accuracy. L1D coverage is calculated per trace as $(M_{\mathrm{baseline}}-M_{\mathrm{prefetcher}})/M_{\mathrm{baseline}}\times100\%$, where $M$ is the count of L1D load misses and the baseline is no-prefetch + LRU.

\subsection{Trace inventory and result provenance}
All reported rows correspond to the three trace names listed above and record 25M warmup plus 25M measured instructions in the raw JSON exports. The CSV and JSON exports are included in this archive. The raw JSON files reference simulator-log paths, but those log files were not included in the data archive supplied for report preparation.

\textbf{Two important provenance limits.} First, the five Task 2B threshold values have identical measured outputs. The supplied \texttt{dyn\_irrip.l2c\_repl} defines \texttt{THRESHOLD 0.25} at compile time and does not read \texttt{CHAMPSIM\_THRESHOLD}; the runner's environment-variable sweep therefore does not actually change this policy. Second, Task 3 result metadata identifies the combined replacement as \texttt{dyn\_irrip}; the separate prefetch-aware implementation is in \texttt{prefetch\_irrip.l2c\_repl}. The supplied Task 3 measurements cannot be claimed as an evaluation of the prefetch-aware file. Both issues are described in their respective sections rather than hidden behind the repeated results.

\newpage\section{Task 1A: Understanding the IP-Stride Prefetcher}
\label{sec:task1a}
\subsection{State and lookup}
The supplied \texttt{prefetcher/}\allowbreak\texttt{ip\_stride.l1d\_pref} implementation defines \texttt{IP\_TRACKER\_COUNT=24}, \texttt{PREFETCH\_DEGREE=3}, and \texttt{MASK\_IP=0xFF}. Each tracker contains \texttt{ip}, \texttt{last\_cl\_addr}, \texttt{last\_stride}, and an LRU rank. The IP is masked to its lower eight bits before lookup, so instruction pointers that differ only above those bits alias to the same logical IP. The implementation does not use a valid bit in its tracker structure.

On an IP miss, it chooses the tracker whose LRU rank is 23, stores the masked IP and current cache-line address, resets the previous stride to zero, updates LRU ranks, and returns without prefetching. On a hit, it calculates a signed stride in cache-line units as current block address minus the previous block address. A zero stride returns immediately. Otherwise, prefetches are issued only when the current stride equals the previously recorded stride. This is a two-consecutive-equal-strides test rather than a saturating confidence counter.

\subsection{Prefetch generation and safety}
The next three block addresses are computed as $A + k\Delta$, for $k=1,2,3$, and converted to byte addresses by shifting by \texttt{LOG2\_BLOCK\_SIZE}. A candidate is discarded, and the loop stops, if it is outside the same 4 KiB page as the current demand access. If MSHR occupancy is below half of MSHR capacity, the prefetch is requested to fill L1; otherwise it is requested to fill L2. The final-statistics callback prints the fixed degree but does not report useful/issued totals.

\subsection{Example state evolution}
The following example uses cache-block numbers (not byte addresses) and assumes every address remains in the same page.
\begin{table}[H]\centering\small
\caption{Example for one masked IP and block sequence 100, 102, 104, 106.}
\begin{tabularx}{\linewidth}{c c c X X}
\toprule Access & Block & Observed stride & Stored state after access & Prefetch action \\
\midrule
1 & 100 & -- & Last block = 100; last stride = 0 & Allocate tracker; none \\
2 & 102 & $+2$ & Last block = 102; last stride = $+2$ & No: previous stride was 0 \\
3 & 104 & $+2$ & Last block = 104; last stride = $+2$ & Yes: blocks 106, 108, 110 \\
4 & 106 & $+2$ & Last block = 106; last stride = $+2$ & Yes: blocks 108, 110, 112 \\
\bottomrule
\end{tabularx}\end{table}

\subsection{Strengths and limitations}
The design is simple and inexpensive, and it can issue multiple requests as soon as a constant stride repeats. It does not encode a multi-stride signature, however, so a repeating sequence such as $+1,+2,+2,+2,+3$ is not modeled as a cycle. It may fail to prefetch across a page boundary even if the stream continues contiguously, and it can lose history when a tracker is evicted. Because tracker IPs are masked and there is no valid bit, a masked IP equal to zero can also match an initially zero-initialized tracker. These limitations motivate the extra history and signature state in Task 1B.

\newpage\section{Task 1B: Stride Pattern Prefetcher (SPP)}
\label{sec:task1b}
\subsection{Data structures and operation}
The supplied \texttt{spp.l1d\_pref} has three direct-mapped structures: 24 IP trackers, a 1,024-entry signature table, and a 512-entry prefetch tracker. The IP key is again masked to eight bits. Each IP tracker stores a 64-bit IP field, last cache-line address, last stride, 32-bit LRU rank, four signed 64-bit stride-history values, a history count, a confidence integer, and a valid flag. When a new IP is seen, an invalid entry is preferred; otherwise an LRU entry is reused and its history is reset.

For each nonzero stride, the implementation updates a short stride-history buffer. After four previous strides are available, it forms a 64-bit signature using the recurrence $s_{i+1}=31s_i+h_i$ over the four signed history values, then uses \texttt{signature \% 1024} as the index. A signature-table entry stores the full signature tag, the next observed stride, a confidence integer, and a valid flag. Matching observations increment confidence up to 7; a changed predicted stride replaces the stored stride and resets confidence to 1. A pattern-table prediction is accepted only when the full signature matches and its confidence is at least 2.

The predictor prioritizes a constant stride whenever its per-IP confidence is at least 1. That confidence is incremented when the current stride equals the previous stride and reset to 0 on a mismatch. If uniform confidence is 0 and four-history state is available, it attempts a signature-table prediction. Thus the source implements a compact conditional fallback, not a separate multi-state classifier with hysteresis. On each demand access it generates up to the current degree of next-stride requests; after each predicted request, the local history is shifted and the next prediction is selected again. A negative candidate or a cross-page candidate terminates the loop. Prefetches fill L1 when MSHR occupancy is below half capacity and L2 otherwise.

\subsection{Accuracy-based degree control}
Accepted prefetch requests increment the global \texttt{issued} counter. A prefetch-table entry, indexed by cache-line address modulo 512, records the address and a state: 1 for an accepted request not yet marked as filled, 2 after a matching prefetch fill, and 0 when cleared. A later demand access that hits a state-2 tracked line increments \texttt{useful}. On each interval where at least 256 additional requests have been accepted, the policy computes the usefulness ratio from the deltas since the last update and sets degree to 1 below 30\%, 2 from 30\% to below 60\%, and 3 at 60\% or above. Degree begins at 3. This is a batch update after issued requests, not a continuously updated sliding window.

\begin{table}[H]\centering\small
\caption{Degree selection encoded in the supplied SPP implementation.}
\begin{tabular}{p{0.45\linewidth}p{0.10\linewidth}p{0.37\linewidth}}
\toprule Measured accuracy since previous update & Degree & Effect \\
\midrule
Below 30\% & 1 & Conservative; limits traffic after low observed usefulness \\
30\% to below 60\% & 2 & Intermediate aggressiveness \\
60\% or above & 3 & Maximum degree in this implementation \\
\bottomrule
\end{tabular}\end{table}

\subsection{Logical storage accounting}
The accounting below counts declared field information bits, treating a boolean as one logical bit and a C++ \texttt{int} as 32 bits. The IP field is counted as 64 bits because the source stores it in a \texttt{uint64\_t}, even though the value is masked to eight bits before use.
\begin{table}[H]\centering\scriptsize
\caption{Bit-level logical storage for all SPP state. C++ object padding is reported separately below.}
\resizebox{\linewidth}{!}{\begin{tabular}{llrrr}
\toprule Structure & Field & Bits per entry & Entries & Total bits \\
\midrule
IP tracker & IP tag storage & 64 & 24 & 1,536 \\
IP tracker & Last cache-line address & 64 & 24 & 1,536 \\
IP tracker & Last stride & 64 & 24 & 1,536 \\
IP tracker & LRU rank & 32 & 24 & 768 \\
IP tracker & History[4] & 256 & 24 & 6,144 \\
IP tracker & History count & 32 & 24 & 768 \\
IP tracker & Confidence & 32 & 24 & 768 \\
IP tracker & Valid flag (logical bit) & 1 & 24 & 24 \\
\midrule
Signature entry & Signature tag & 64 & 1,024 & 65,536 \\
Signature entry & Predicted stride & 64 & 1,024 & 65,536 \\
Signature entry & Confidence & 32 & 1,024 & 32,768 \\
Signature entry & Valid flag (logical bit) & 1 & 1,024 & 1,024 \\
\midrule
Prefetch tracker & Cache-line address & 64 & 512 & 32,768 \\
Prefetch tracker & State integer & 32 & 512 & 16,384 \\
\midrule
Global state & issued, useful, last\_issued, last\_useful & 256 & 1 & 256 \\
Global state & Degree integer & 32 & 1 & 32 \\
\midrule
\multicolumn{4}{r}{\textbf{Total logical field bits}} & \textbf{227,384} \\
\multicolumn{4}{r}{\textbf{Logical bytes (bits/8)}} & \textbf{28,423} \\
\multicolumn{4}{r}{\textbf{Logical KiB (bytes/1024)}} & \textbf{27.76} \\
\bottomrule\end{tabular}}\end{table}

On a conventional 64-bit ABI, typical structure alignment would make the arrays approximately $24\times80 + 1024\times24 + 512\times16=34,688$ bytes. Adding the four 64-bit counters and the 32-bit degree gives about 34,724 bytes (33.91 KiB), excluding executable/runtime overhead. The logical-field total and ABI-size estimate answer different questions: the former counts information field widths, while the latter includes expected padding. Exact physical bytes should be confirmed with \texttt{sizeof} on the compiler used for the final submission.

\subsection{Measured results}
\begin{table}[H]\centering\scriptsize
\caption{Task 1B measurements. Coverage and accuracy are not applicable to the no-prefetch baseline.}
\resizebox{\linewidth}{!}{\begin{tabular}{llrrrrr}
\toprule Trace & Configuration & IPC & Speedup & L1D Load MPKI & Accuracy (\%) & Coverage (\%) \\
\midrule
Tr1.gz & No prefetch + LRU & 0.232345 & 1.000000 & 10.6780 & N/A & N/A \\
Tr1.gz & IP-stride & 0.233452 & 1.004764 & 10.1190 & 52.71 & 5.24 \\
Tr1.gz & SPP & 0.232577 & 1.000999 & 10.5465 & 41.33 & 1.23 \\
Tr2.gz & No prefetch + LRU & 0.251158 & 1.000000 & 13.5379 & N/A & N/A \\
Tr2.gz & IP-stride & 0.251260 & 1.000406 & 13.4378 & 38.18 & 0.74 \\
Tr2.gz & SPP & 0.251094 & 0.999745 & 13.5014 & 33.54 & 0.27 \\
Tr3.gz & No prefetch + LRU & 0.382065 & 1.000000 & 3.4034 & N/A & N/A \\
Tr3.gz & IP-stride & 0.382862 & 1.002086 & 3.3182 & 15.99 & 2.51 \\
Tr3.gz & SPP & 0.382139 & 1.000194 & 3.3938 & 8.51 & 0.28 \\
\bottomrule\end{tabular}}
\end{table}

\plotpair{plots/task1b/ipc.png}{IPC}{plots/task1b/speedup.png}{Speedup}{Task 1B IPC and speedup across traces}{fig:t1b-ipc-speedup}
\plotpair{plots/task1b/l1d_load_mpki.png}{L1D load MPKI}{plots/task1b/prefetch_accuracy_pct.png}{Prefetch accuracy}{Task 1B miss intensity and prefetch accuracy}{fig:t1b-mpki-accuracy}
\plotpair{plots/task1b/l1d_coverage_pct.png}{L1D coverage}{plots/task1b/ipc_change_vs_baseline_pct.png}{IPC change vs. baseline}{Task 1B coverage and relative IPC change}{fig:t1b-coverage-change}
\plotone{plots/task1b/l1d_mpki_reduction_vs_baseline_pct.png}{L1D load-MPKI reduction relative to no-prefetch + LRU. Positive values mean fewer L1D load misses.}{fig:t1b-mpki-reduction}

\subsection{Interpretation}
On Tr1, IP-stride improves IPC by 0.476\% over no-prefetch + LRU, while SPP improves it by 0.100\%. On Tr2, IP-stride improves by 0.041\%, while SPP is 0.025\% below baseline. On Tr3, IP-stride improves by 0.209\% and SPP by 0.019\%. Thus IP-stride has the highest IPC of the compared configurations on all three traces in this data set.

IP-stride also has higher measured accuracy and coverage than SPP on every trace. The L1D load-MPKI reductions are 5.24\%, 0.74\%, and 2.51\% for IP-stride on Tr1/Tr2/Tr3, compared with 1.23\%, 0.27\%, and 0.28\% for SPP. These observations do not prove that multi-stride prediction is generally worse; they show that the current compact signature/fallback implementation does not outperform the provided IP-stride baseline on these particular traces. The signature table's direct mapping, small IP tracker pool, masked IP key, confidence fallback, and modulo-indexed prefetch tracker are plausible limits, but the measurements alone do not isolate which mechanism explains the gap.

\newpage% ============================================================
% TASK 2A: Instruction-Conscious L2 Replacement Policy (I-RRIP)
% ============================================================
% TODO Task 2A: implement, measure, plot, and explain I-RRIP.
\section{Task 2A: Instruction-Conscious L2 Replacement Policy (I-RRIP)}

\subsection{Baseline RRIP and Motivation}
\TODO{Briefly describe the provided RRIP policy and why prioritizing instruction lines at L2 may matter for the evaluated workloads. Reference the repository's actual RRIP implementation.}

\subsection{Implementation Changes}
The required changes cover replacement-state update/insertion, promotion, and eviction, with the policy informed by whether a packet represents an instruction line (the assignment notes \texttt{packet.instruction}).
\begin{itemize}
    \item \textbf{API / metadata changes:} \TODO{Describe the signature change to update replacement state and all affected call sites.}
    \item \textbf{Insertion:} \TODO{Specify initial RRPV/state for instruction and data lines and explain the rationale.}
    \item \textbf{Promotion:} \TODO{Specify hit behavior for instruction and data lines. Explain whether the classes promote differently.}
    \item \textbf{Eviction:} \TODO{Explain victim selection and how ties, aging, and saturation are handled.}
    \item \textbf{Compatibility:} \TODO{Describe API, compilation, and corner-case checks performed.}
\end{itemize}
\begin{figure}[H]\centering
\evidencebox[4.2cm]{\TODO{Insert a compact policy diagram or pseudocode comparing RRIP and I-RRIP for insertion, hit promotion, aging, and victim selection.}}
\caption{Instruction-conscious modifications to the RRIP lifecycle.}\label{fig:irrip_policy}
\end{figure}

\subsection{Experimental Results}
Report IPC, speedup, and load MPKI for L2C, L2 instruction, and L2 data. Include baseline RRIP and I-RRIP results for every trace.
\begin{table}[H]\centering
\caption{Task 2A: I-RRIP result table. MPKI columns must specify the exact ChampSim statistic used.}
\small
\resizebox{\linewidth}{!}{%
\begin{tabular}{llrrrrr}
\toprule Trace & Policy & IPC & Speedup & L2C Load MPKI & L2I Load MPKI & L2D Load MPKI \\
\midrule
\TODO{trace} & RRIP baseline & \TODO{value} & 1.00$\times$ & \TODO{value} & \TODO{value} & \TODO{value} \\
\TODO{trace} & I-RRIP & \TODO{value} & \TODO{value} & \TODO{value} & \TODO{value} & \TODO{value} \\
\TODO{add rows} & & & & & & \\
\bottomrule
\end{tabular}%
}\end{table}

\begin{figure}[H]\centering
\begin{subfigure}{0.48\textwidth}\centering\smallplotplaceholder{IPC per trace: RRIP baseline vs. I-RRIP}\caption{IPC}\end{subfigure}\hfill
\begin{subfigure}{0.48\textwidth}\centering\smallplotplaceholder{Speedup per trace against RRIP baseline}\caption{Speedup}\end{subfigure}
\vspace{0.6em}
\begin{subfigure}{0.48\textwidth}\centering\smallplotplaceholder{L2C load MPKI per trace and policy}\caption{L2C load MPKI}\end{subfigure}\hfill
\begin{subfigure}{0.48\textwidth}\centering\smallplotplaceholder{L2 instruction-load MPKI per trace and policy}\caption{L2 instruction MPKI}\end{subfigure}
\vspace{0.6em}
\begin{subfigure}{0.48\textwidth}\centering\smallplotplaceholder{L2 data-load MPKI per trace and policy}\caption{L2 data MPKI}\end{subfigure}
\caption{Required Task 2A performance metrics.}\label{fig:irrip_results}
\end{figure}

\subsection{Analysis and Rationale}
\begin{enumerate}
    \item \textbf{Insertion logic:} \TODO{Explain why the implemented insertion priorities preserve or demote instruction/data lines.}
    \item \textbf{Promotion logic:} \TODO{Discuss whether instruction hits receive stronger protection and the effect on data lines.}
    \item \textbf{Eviction logic:} \TODO{Explain victim selection and how instruction and data RRPVs compete.}
    \item \textbf{Results:} \TODO{Explain changes in IPC and all three MPKI metrics, including positive and negative cases. Avoid claiming causality beyond measured evidence.}
\end{enumerate}

\newpage\section{Task 2B: Dynamic I-RRIP with Set Dueling}
\label{sec:task2b}
\subsection{Set selection and decision statistic}
The supplied \texttt{dyn\_irrip.l2c\_repl} sets \texttt{maxRRPV=3}, uses \texttt{RESET\_INTERVAL=4096}, and defines the threshold as the compile-time constant \texttt{THRESHOLD 0.25}. For set index $s$, sets with $s\bmod32=0$ are RRIP leader sets, sets with $s\bmod32=1$ are I-RRIP leader sets, and all other sets use the global follower decision. The two leader classes together account for 2 of every 32 sets (6.25\% when the number of sets is divisible by 32), with each leader class representing 3.125\% of sets.

The policy updates shared counters on accesses to either leader class and computes
\[
R=\frac{A_I+M_I}{A_{\mathrm{total}}+M_{\mathrm{total}}},
\]
where the \texttt{instr\_accesses}, \texttt{instr\_misses}, \texttt{total\_accesses}, and \texttt{total\_misses} counters implement these quantities. The ratio is recalculated after each leader-set callback and compared with 0.25; follower sets use I-RRIP when $R\geq0.25$ and RRIP otherwise. Once total accesses reach 4,096, the counters are reset after the ratio update. The source uses a single shared statistic rather than keeping independent utility counters for the two leader policies, so its decision is instruction-activity based rather than a direct comparison of measured miss cost between the two leader groups.

\subsection{Threshold sensitivity: a source-level issue}
The runner attempted thresholds 0.10, 0.20, 0.25, 0.30, and 0.40. The measured IPC values and cache metrics repeat exactly across those five rows. This is explained by the implementation: \texttt{dyn\_irrip.l2c\_repl} hardcodes \texttt{THRESHOLD 0.25} and never reads the \texttt{CHAMPSIM\_THRESHOLD} environment variable that the runner changes. Therefore these runs all use the same effective threshold and do \emph{not} constitute a valid threshold sensitivity experiment. The repeated table below preserves the observed results for reproducibility; it should not be used to claim that the policy is insensitive to threshold or that 0.25 is optimal.

\begin{table}[H]\centering\small
\caption{Observed Dynamic I-RRIP IPC by requested threshold. The invariance is an artifact of the compile-time constant, not evidence from a working sweep.}
\begin{tabular}{lrrrr}
\toprule Requested threshold & Tr1 IPC & Tr2 IPC & Tr3 IPC & Mean speedup vs RRIP \\
\midrule
0.10 & 0.232281 & 0.256285 & 0.381742 & 1.006004 \\
0.20 & 0.232281 & 0.256285 & 0.381742 & 1.006004 \\
0.25 & 0.232281 & 0.256285 & 0.381742 & 1.006004 \\
0.30 & 0.232281 & 0.256285 & 0.381742 & 1.006004 \\
0.40 & 0.232281 & 0.256285 & 0.381742 & 1.006004 \\
\bottomrule\end{tabular}\end{table}
\plotpair{plots/task2b/sensitivity_ipc.png}{IPC by requested threshold}{plots/task2b/sensitivity_l2i_load_mpki.png}{L2 instruction MPKI by requested threshold}{Requested-threshold sweep; identical bars reflect the hardcoded threshold described above}{fig:t2b-sensitivity}

\subsection{Measured policy comparison at effective threshold 0.25}
\begin{table}[H]\centering\scriptsize
\caption{Task 2B results. Speedup is relative to RRIP on the same trace.}
\resizebox{\linewidth}{!}{\begin{tabular}{llrrrrr}
\toprule Trace & Policy & IPC & Speedup & L2C Load MPKI & L2 Instruction MPKI & L2 Data MPKI \\
\midrule
Tr1.gz & RRIP baseline & 0.232231 & 1.000000 & 2.2021 & 10.5618 & 7.0769 \\
Tr1.gz & I-RRIP & 0.232164 & 0.999711 & 2.4089 & 9.8916 & 7.7172 \\
Tr1.gz & Dynamic I-RRIP (compiled threshold 0.25) & 0.232281 & 1.000215 & 2.3941 & 9.9094 & 7.6956 \\
Tr2.gz & RRIP baseline & 0.251806 & 1.000000 & 2.1711 & 17.7813 & 7.9645 \\
Tr2.gz & I-RRIP & 0.256625 & 1.019138 & 2.3754 & 13.0364 & 9.1737 \\
Tr2.gz & Dynamic I-RRIP (compiled threshold 0.25) & 0.256285 & 1.017788 & 2.3631 & 13.3453 & 9.1221 \\
Tr3.gz & RRIP baseline & 0.381739 & 1.000000 & 0.7090 & 1.1970 & 2.8424 \\
Tr3.gz & I-RRIP & 0.381981 & 1.000634 & 0.7264 & 1.0221 & 2.8761 \\
Tr3.gz & Dynamic I-RRIP (compiled threshold 0.25) & 0.381742 & 1.000008 & 0.7235 & 1.1464 & 2.8764 \\
\bottomrule\end{tabular}}
\end{table}
\plotpair{plots/task2b/ipc.png}{IPC}{plots/task2b/speedup.png}{Speedup}{Task 2B policy-level IPC and speedup}{fig:t2b-ipc-speedup}
\plotpair{plots/task2b/l2c_load_mpki.png}{L2C load MPKI}{plots/task2b/l2i_load_mpki.png}{L2 instruction-load MPKI}{Task 2B aggregate and instruction MPKI}{fig:t2b-l2c-l2i}
\plotpair{plots/task2b/l2d_load_mpki.png}{L2 data-load MPKI}{plots/task2b/ipc_change_vs_rrip_by_threshold_pct.png}{IPC change vs. RRIP (\%)}{Task 2B data-MPKI trade-off and relative IPC change}{fig:t2b-l2d-change}
\plotone{plots/task2b/l2i_mpki_reduction_vs_rrip_by_threshold_pct.png}{Instruction-load MPKI reduction relative to RRIP for the dynamic policy at the reported threshold.}{fig:t2b-l2i-reduction}

\subsection{Interpretation and required correction}
At the measured effective threshold of 0.25, dynamic I-RRIP changes IPC by +0.0215\% on Tr1, +1.7788\% on Tr2, and approximately +0.0008\% on Tr3 relative to RRIP. L2 instruction MPKI declines on all three traces, while aggregate L2C and data MPKI rise. Tr2 is the clearest performance win in the supplied data; Tr1 and Tr3 are essentially neutral in IPC.

To complete the assignment's sensitivity requirement, replace the compile-time macro with a runtime parameter read from \texttt{CHAMPSIM\_THRESHOLD} (with a safe default and range check), rebuild for every threshold or ensure the runtime argument reaches the policy, and rerun all three traces at the five values. Keep leader-set placement and all other parameters fixed. Until those measurements exist, no threshold selection rationale can be inferred from this sweep.

\newpage\section{Task 3: Combining SPP with Dynamic I-RRIP}
\label{sec:task3}
\subsection{Prefetch-aware replacement code in the supplied archive}
The code archive contains a separate \texttt{replacement/prefetch\_irrip.l2c\_repl} implementation. That file tracks a per-set/per-way \texttt{prefetched} flag, uses the \texttt{PREFETCH} request type, and reads \texttt{CHAMPSIM\_THRESHOLD} at initialization (default 0.25). When a prefetch fills a way, the flag is set and its RRPV is initialized to the maximum value 3, leaving it easy to evict if it is never used. During victim search it first prefers a maximum-RRPV prefetched line; if none qualifies and I-RRIP is active, it next prefers a maximum-RRPV data line, then falls back to any maximum-RRPV line. When a demand hit occurs on a flagged prefetched line, the code counts it as useful, clears the flag, and promotes it using the active RRIP/I-RRIP rules. A prefetch hit itself returns without promoting the line. The policy excludes prefetch requests from the counters used to choose the dynamic mode.

This design follows a reasonable pollution-control rule: speculative lines enter with low protection, while demand reuse gives them normal replacement treatment. The implementation also reports useful and unused-prefetch counts at finalization. These are code-review facts about \texttt{prefetch\_irrip}; they are not, by themselves, evidence of performance improvement.

\subsection{Which Task 3 configuration was measured?}
The Task 3 CSV/JSON metadata and run labels identify the combined
configuration as \texttt{combined-dyn\_irrip}. The supplied runner also
defaults the combined replacement option to \texttt{dyn\_irrip}. In the
provided source archive, \texttt{dyn\_irrip.l2c\_repl} does not implement
the per-way prefetch tracking/eviction logic described above. Consequently,
the measurements below are labeled \textbf{SPP + dyn\_irrip (as recorded)};
they cannot be presented as measured results for the separate prefetch-aware
\texttt{prefetch\_irrip} file. The Task 3 requirement should be re-evaluated
after the prefetch-aware logic is integrated into the specified dynamic policy
and the runs are repeated.

\subsection{Measured results: IPC, speedup, and L1D behavior}
\begin{table}[H]\centering\scriptsize
\caption{Task 3 data as recorded. Speedup is relative to no-prefetch + LRU for the same trace.}
\resizebox{\linewidth}{!}{\begin{tabular}{llrrrrr}
\toprule Trace & Configuration & IPC & Speedup & Accuracy (\%) & Coverage (\%) & L1D load MPKI \\
\midrule
Tr1.gz & Baseline (LRU, no prefetch) & 0.232345 & 1.000000 & N/A & N/A & 10.6780 \\
Tr1.gz & SPP only & 0.232577 & 1.000999 & 41.33 & 1.23 & 10.5465 \\
Tr1.gz & Dynamic I-RRIP only & 0.232281 & 0.999725 & N/A & N/A & 10.6788 \\
Tr1.gz & SPP + \texttt{dyn\_irrip} (as recorded) & 0.232482 & 1.000590 & 41.07 & 1.22 & 10.5476 \\
Tr2.gz & Baseline (LRU, no prefetch) & 0.251158 & 1.000000 & N/A & N/A & 13.5379 \\
Tr2.gz & SPP only & 0.251094 & 0.999745 & 33.54 & 0.27 & 13.5014 \\
Tr2.gz & Dynamic I-RRIP only & 0.256285 & 1.020413 & N/A & N/A & 13.5412 \\
Tr2.gz & SPP + \texttt{dyn\_irrip} (as recorded) & 0.256303 & 1.020485 & 33.34 & 0.23 & 13.5064 \\
Tr3.gz & Baseline (LRU, no prefetch) & 0.382065 & 1.000000 & N/A & N/A & 3.4034 \\
Tr3.gz & SPP only & 0.382139 & 1.000194 & 8.51 & 0.28 & 3.3938 \\
Tr3.gz & Dynamic I-RRIP only & 0.381742 & 0.999155 & N/A & N/A & 3.4033 \\
Tr3.gz & SPP + \texttt{dyn\_irrip} (as recorded) & 0.381886 & 0.999531 & 8.52 & 0.29 & 3.3937 \\
\bottomrule\end{tabular}}
\end{table}

\subsection{L2 behavior}
\begin{table}[H]\centering\scriptsize
\caption{L2 statistics by trace/configuration. These use the separate aggregate L2C and instruction/data counters described in Section 1.}
\resizebox{\linewidth}{!}{\begin{tabular}{llrrr}
\toprule Trace & Configuration & L2C load MPKI & L2 instruction-load MPKI & L2 data-load MPKI \\
\midrule
Tr1.gz & Baseline (LRU, no prefetch) & 2.2774 & 10.3377 & 7.3580 \\
Tr1.gz & SPP only & 2.2805 & 10.3556 & 7.2210 \\
Tr1.gz & Dynamic I-RRIP only & 2.3941 & 9.9094 & 7.6956 \\
Tr1.gz & SPP + \texttt{dyn\_irrip} (as recorded) & 2.3894 & 9.9151 & 7.5314 \\
Tr2.gz & Baseline (LRU, no prefetch) & 2.2437 & 18.8533 & 8.5035 \\
Tr2.gz & SPP only & 2.2457 & 18.9067 & 8.4938 \\
Tr2.gz & Dynamic I-RRIP only & 2.3631 & 13.3453 & 9.1221 \\
Tr2.gz & SPP + \texttt{dyn\_irrip} (as recorded) & 2.3648 & 13.3179 & 9.1141 \\
Tr3.gz & Baseline (LRU, no prefetch) & 0.6972 & 1.1241 & 2.8302 \\
Tr3.gz & SPP only & 0.7033 & 1.1386 & 2.8119 \\
Tr3.gz & Dynamic I-RRIP only & 0.7235 & 1.1464 & 2.8764 \\
Tr3.gz & SPP + \texttt{dyn\_irrip} (as recorded) & 0.7275 & 1.1494 & 2.8502 \\
\bottomrule\end{tabular}}
\end{table}

\plotpair{plots/task3/ipc.png}{IPC}{plots/task3/speedup.png}{Speedup}{Task 3 measured configuration comparison}{fig:t3-ipc-speedup}
\plotpair{plots/task3/l1d_load_mpki.png}{L1D load MPKI}{plots/task3/l1d_coverage_pct.png}{L1D coverage}{Task 3 L1D miss and coverage behavior}{fig:t3-l1d}
\plotpair{plots/task3/prefetch_accuracy_pct.png}{Prefetch accuracy}{plots/task3/l2c_load_mpki.png}{L2C load MPKI}{Task 3 prefetch accuracy and aggregate L2C metric}{fig:t3-accuracy-l2c}
\plotpair{plots/task3/l2i_load_mpki.png}{L2 instruction-load MPKI}{plots/task3/l2d_load_mpki.png}{L2 data-load MPKI}{Task 3 L2 instruction and data behavior}{fig:t3-l2i-l2d}
\plotpair{plots/task3/ipc_change_vs_baseline_pct.png}{IPC change vs. baseline (\%)}{plots/task3/ipc_gain_vs_best_standalone_pct.png}{IPC gain vs. best standalone (\%)}{Additional Task 3 relative-performance views}{fig:t3-relative}

\subsection{Does the measured combined configuration meet the IPC goal?}
The answer is trace-dependent even for the configuration recorded in the CSVs. On Tr1, the combined IPC is 0.232482, above the baseline and dynamic-only result but below SPP-only (0.232577). On Tr2, combined IPC is 0.256303, just above dynamic-only (0.256285) and clearly above SPP-only (0.251094); the difference over dynamic-only is only about 0.007\%. On Tr3, combined IPC is 0.381886, above dynamic-only (0.381742) but below SPP-only (0.382139) and below the no-prefetch baseline (0.382065). Therefore the recorded configuration does not outperform both standalone methods on all traces.

The results also show why no single metric is enough. Prefetch accuracy is approximately 41.07\%, 33.34\%, and 8.52\% for the combined configuration on Tr1--Tr3, while coverage is about 1.22\%, 0.23\%, and 0.29\%. L1D MPKI is modestly lower than baseline on all traces, but aggregate L2C MPKI is higher on all three. At the same time, instruction-load MPKI is considerably lower on Tr2, while data MPKI is higher than the no-prefetch baseline on all traces. Because the combined runs used \texttt{dyn\_irrip}, these observations are about the measured SPP + dynamic replacement configuration and must not be used to infer the effect of prefetch-aware victim selection.

\newpage
\section{Overall Comparison and Conclusion}
\subsection{Cross-task findings}
\begin{itemize}
\item \textbf{Task 1B:} IP-stride has the highest IPC of the tested prefetcher configurations on all three traces. SPP reduces L1D MPKI relative to the no-prefetch baseline on all traces, but less than IP-stride, and has lower reported accuracy than IP-stride on all three traces.
\item \textbf{Task 2A:} I-RRIP lowers L2 instruction-load MPKI on all three traces, but increases L2 data-load MPKI and aggregate L2C load MPKI. IPC improves most on Tr2; Tr1 is marginally lower than the RRIP baseline.
\item \textbf{Task 2B:} At the compiled threshold of 0.25, dynamic I-RRIP gives its largest IPC gain on Tr2. The supplied five-threshold run is not a valid sensitivity study because the policy has a compile-time threshold; no optimal threshold can be selected from these data.
\item \textbf{Task 3:} The combined result as recorded is better than standalone dynamic I-RRIP by a very small margin on Tr2, but does not beat SPP on Tr1 or Tr3. Because the experiment used \texttt{dyn\_irrip} rather than the separate prefetch-aware implementation, this result does not establish whether prefetch-aware replacement meets the assignment's goal.
\end{itemize}

\subsection{Limitations and next steps}
The most important follow-up is to make the dynamic threshold truly runtime-configurable, rebuild, and rerun the five-threshold sweep. For Task 3, integrate the prefetch-awareness behavior into the policy actually named by the runner (the assignment specifies \texttt{dyn\_irrip.l2c\_repl}) or run the supplied \texttt{prefetch\_irrip} replacement explicitly; then regenerate the result CSVs and plots from those new logs. If small IPC differences are to be emphasized, repeat runs or another measure of run-to-run variability would also help. Trace semantic descriptions, compiler/host metadata, and complete raw logs should be retained for a reproducible final submission.

\subsection{Conclusion}
The measured results demonstrate that instruction-conscious replacement can reduce instruction MPKI, but at the cost of higher data MPKI on these traces. The supplied SPP achieves only small coverage gains, and IP-stride outperforms it in the current data. Most importantly, the current artifacts expose two experiment-validity gaps—the ineffective Task 2B sweep and a Task 3 policy/configuration mismatch—which should be corrected before asserting that the complete assignment goals have been met.
\end{document}
